% LedModule
\section{【基本】LedModule — 最小構成の基本パターン}

\begin{patternbox}{LedModule で学べるパターン}
\begin{itemize}
    \item Config構造体によるピンアサインの外部注入
    \item Data構造体のデフォルト値
    \item .h/.cpp分離（循環依存回避）
    \item 出力モジュールの基本形（SystemDataの値に従って動作するだけ）
\end{itemize}
\end{patternbox}

\subsection{ヘッダー（LedModule.h）}

\begin{lstlisting}[style=cppstyle, caption={LedModule.h}]
#pragma once
#include <Arduino.h>
#include "IModule.h"

struct LedConfig {
    uint8_t ledPin;  // LED出力ピン
};

struct LedData {
    bool state = false;       // LED状態（デフォルト: 消灯）
    uint8_t brightness = 0;   // 輝度（将来拡張用）
};

class LedModule : public IModule {
private:
    LedConfig _config;
public:
    LedModule(const LedConfig& config);
    bool init() override;
    void update(SystemData& data) override;
};
\end{lstlisting}

\begin{pointbox}[Config構造体のポイント]
\texttt{LedConfig}にはピン番号だけを持つ。
ピン番号を\texttt{LedModule}内にハードコードせず、Config構造体で外部から注入する。
これにより、同じ\texttt{LedModule}クラスを別プロジェクトでピンだけ変えて再利用できる。
\end{pointbox}

\begin{pointbox}[Data構造体のデフォルト値]
\texttt{LedData}の\texttt{state = false}は「\texttt{update()}が1度も呼ばれていない状態でも安全」を保証する。
ロジックフェーズが\texttt{state}を\texttt{true}に変更するまで、LEDは消灯したままとなる。
\end{pointbox}

\subsection{実装（LedModule.cpp）}

\begin{lstlisting}[style=cppstyle, caption={LedModule.cpp}]
#include "LedModule.h"
#include "SystemData.h"  // ヘッダーでは前方宣言のみ、cppで完全な定義を取得

LedModule::LedModule(const LedConfig& config) : _config(config) {}

bool LedModule::init() {
    pinMode(_config.ledPin, OUTPUT);
    Serial.println("[Led] init OK");
    return true;
}

void LedModule::update(SystemData& data) {
    digitalWrite(_config.ledPin, data.led.state ? HIGH : LOW);
}
\end{lstlisting}

\begin{pointbox}[出力モジュールの責務]
\texttt{update()}は\textbf{自身のData構造体の値に従って動作するだけ}。
「いつ点灯するか」「何の条件で消すか」などのロジックは一切書かない。
それはロジックフェーズ（\texttt{applyPattern()}等）の仕事である。
\end{pointbox}

\begin{pointbox}[.h/.cpp分離のルール]
ヘッダーでは\texttt{SystemData.h}をインクルードせず、\texttt{IModule.h}内の前方宣言だけを利用する。
\texttt{SystemData.h}のインクルードは\texttt{.cpp}側で行う。
これにより\texttt{LedModule.h} → \texttt{SystemData} → \texttt{LedModule.h}の循環依存を防ぐ。
\end{pointbox}

\subsection{ProjectConfig.hでの定義}

\begin{lstlisting}[style=cppstyle, caption={ProjectConfig.hでのConfig定義}]
const LedConfig LED_CONFIG = {
    .ledPin = 2,  // GPIO2: オンボードLED
};
\end{lstlisting}

回路設計者はこのファイルだけ見ればピンアサインが分かる。
\texttt{LED\_CONFIG}という大文字スネークケースの命名規則も守っている。
